\documentclass[12pt]{article}
\usepackage[utf8]{inputenc}
\usepackage[T1]{fontenc}
\usepackage{hyperref}
\title{Pai Tavytera Language Resources for K-12 Curriculum Development}
\author{Compiled by Claude Sonnet 5.5 (claude-sonnet-5-5) for Andr\'as Kornai}
\date{}
\begin{document}
\maketitle

\section{Introduction}
Pai Tavytera (Pa\~i Tavyter\~a; ISO 639-3 \texttt{pta}; Glottocode pait1247) is a Tupi-Guarani language of eastern Paraguay (Amambay and Canindey\'u), with roughly 600 speakers by a 2007 estimate. It is written in Latin script, and shift toward Paraguayan Guaran\'i is reported. Documentation is mostly academic and community-led and is largely in Spanish; no complete school curriculum was found.

\section{Primary Educational Resources}
\begin{itemize}
\item Universidad Nacional de Concepci\'on, intercultural education final report (in Spanish), documenting the social and cultural situation of the Pa\~i Tavyter\~a in Amambay, written with Pa\~i Tavyter\~a students; relevant to intercultural teacher training. \url{https://www.fhyce.edu.py/wp-content/uploads/2023/10/INTERCULTURALIDAD-INFORME-FINAL-PARA-DEFENSA.pdf}
\end{itemize}

\section{Linguistic and Reference Materials}
\begin{itemize}
\item Celeste Escobar, \emph{Vocabulario Pai Tavytera Guarani}: vocabulary article based on a corpus of 621 lexemes, in Spanish with an English abstract. \url{https://repositorio.uni.edu.py/documentos/Celeste%20Escobar_Vocabulario%20Pai%20Tavytera%20Guarani.pdf}
\item Wikipedia, ``Pai Tavytera language'': overview, writing system and relationship to Ka'iwa. \url{https://en.wikipedia.org/wiki/Pai_Tavytera_language}
\item Endangered Languages Project entry (Spanish interface) for Pai Tavytera, which describes the Tape Ayvu platform. \url{https://endangeredlanguages.com/es/node/114906}
\end{itemize}

\section{Religious and Community Resources}
No scripture translations were found in search results. Native Land's language listing may help locate community context. \url{https://native-land.ca/listings/languages/pai-tavytera}

\section{Digital and Audio Resources}
According to the Endangered Languages Project entry, Tape Ayvu (launched 2023 by Escobar with community members) offers vocabulary, oral texts and learner materials in bilingual form. I did not locate the platform's own URL in search results.

\section{Curriculum Development Context}
The language is spoken in every generation but is not known to be taught in schools. A curriculum would need to build on Paraguay's intercultural bilingual education framework, community-authored oral texts and the vocabulary work above, with Pa\~i Tavyter\~a teachers leading. Contacting Paraguay's language policy secretariat (Secretar\'ia de Pol\'iticas Ling\"u\'isticas) was suggested by a source but not verified.

\end{document}